% AUTO-GENERATED by figs/make_figures.py -- do not edit.
\begin{figure}[t]\centering
\begin{tikzpicture}
\begin{axis}[paperfig,width=\columnwidth,height=0.6\columnwidth,
  xlabel={gap coefficient of variation $\mathrm{CV}(V)$},
  ylabel={age (normalized)},legend pos=north west,grid=both,
  legend cell align=left]
\addplot[cB,only marks,mark=square*,mark size=1.7pt] coordinates {(0.0,0.41) (0.3,0.44625348238863416) (0.5,0.5108772600771613) (1.0,0.8149996897795648) (1.5,1.3224004704089134) (2.0,2.030147837685105)};
\addlegendentry{mean AoI}
\addplot[cA,only marks,mark=*,mark size=1.7pt] coordinates {(0.0,0.9049999999999999) (0.3,0.9046423617422175) (0.5,0.9044223082534586) (1.0,0.9043348492407165) (1.5,0.9039569330472569) (2.0,0.9033825094312596)};
\addlegendentry{mean OPAoI}
\draw[gray,dashed] (axis cs:1.106,\pgfkeysvalueof{/pgfplots/ymin}) -- (axis cs:1.106,\pgfkeysvalueof{/pgfplots/ymax});
\end{axis}
\end{tikzpicture}
\caption{The AoI/OPAoI inversion (E12). Mean OPAoI ($=T_s+
\mathbb{E}[V]$) is flat in the gap variability, while mean AoI grows
with it; they cross at the exact threshold
$\mathrm{CV}^2(V)=1+2\mathbb{E}[U]/\mathbb{E}[V]$ (dashed:
$\mathrm{CV}\approx1.11$ for the plotted $\mathbb{E}[U]/\mathbb{E}[V]=1/9$),
so for high-variance
terrestrial contacts the outage peak sits {\itshape below}
average AoI.}
\label{fig:inversion}
\end{figure}
